\chapter*{Abstract}
\cftaddtitleline{toc}{chapter}{ABSTRACT}{\thepage}

\emph{Whether a corroded steel component remains fit for service depends on
how deep the metal loss goes as well as how much area it covers. Camera-based
deep learning detects corrosion without contact, but a color camera records
appearance only: it confuses stains, coatings and shadows with corrosion, cannot
see beneath paint, and reports severity as area. Thermal imaging responds to
defects below the surface and LiDAR measures geometry, yet of the 33
corrosion-sensing studies reviewed, 16 rely on RGB alone, one adds thermal
imaging and none fuses RGB, thermal and LiDAR data. This study aims to develop and
evaluate a framework that detects corroded regions by fusing the three sensors,
measures each region's area and metal loss from the LiDAR, and determines what
each sensor contributes. The sensors, carried on an unmanned aerial vehicle, will
be registered to a common frame and stacked into one five-channel input to a
pretrained YOLO11 detector. An ablation will compare single, paired and full
configurations, with a blank-channel control, under identical conditions and a
decision rule fixed in advance. A preliminary test on MID-1K, a public
RGB--thermal--LiDAR person-detection dataset, ran end to end across six
configurations and three seeds. Early fusion reached a test mAP@0.5:0.95 of 0.324
against 0.311 for RGB, within the seed spread and so short of the decision rule,
and thermal fell from 0.756 to 0.521 mAP@0.5 when scored against RGB-aligned
labels, suggesting that registration must precede fusion. The expected
contributions are a registered trimodal corrosion dataset and a controlled
measure of which sensors are worth deploying.}

\vspace{0.8\baselineskip}

\noindent\textbf{Keywords:} \emph{Corrosion detection; multimodal sensor
fusion; early fusion; object detection; YOLO; thermal imaging; LiDAR; ablation
study.}
